\section{Game Logic}\label{sec:game}
The game simulation is deterministic (fixed seeds, no wall-clock time), CPU-only and covered by tests T100--T113. This
section covers what happens to an intruder and to the projectiles; the two-UAV engagement, waves and score are in
Section~\ref{sec:squad}, and the rover's hiding and infiltration planner (MicroMouse flood fill and line follower) in
Section~\ref{sec:rover}. The code is in \code{Scene.cpp} (life cycle, blasts, debris) and \code{Combat.cpp} (weapons, explosions, particles).

\subsection{Intruder life cycle}
\autofig{diagrams}{G3_game_loop}
Each intruder rover is a slot with a \code{Life} state (Fig.~\ref{fig:life}). The mission (Section~\ref{sec:squad}) fills a free slot
with a rover at the compound edge, facing a camp asset; a rover that reaches the asset is a breach and its slot is freed at once.
\begin{figure}[h]
\centering
\begin{tikzpicture}[node distance=14mm and 18mm, every node/.style={font=\scriptsize},
  st/.style={draw,rounded corners=2pt,minimum height=7mm,minimum width=18mm,align=center,fill=blue!5},
  arr/.style={-{Latex[length=1.6mm]},thick}]
  \node[st] (w) {Waiting\\slot free};
  \node[st,right=of w] (a) {Alive\\infiltrates};
  \node[st,right=of a] (r) {Wreck\\burns 15\,s, removed 25\,s};
  \node[st,right=of r] (g) {Gone\\3\,s};
  \draw[arr] (w) -- node[above]{spawn} (a);
  \draw[arr] (a) -- node[above]{blast} (r);
  \draw[arr] (r) -- node[above]{25\,s} (g);
  \draw[arr] (g.south) -- ++(0,-8mm) -| node[pos=0.25,below]{slot free for the next wave} (w.south);
\end{tikzpicture}
\caption{Intruder life cycle (\code{Intruder::Life}); constants \code{kWreckBurnTime}=15\,s, \code{kWreckClearTime}=25\,s,
\code{kRespawnDelay}=3\,s; a breach returns the slot directly to Waiting.}
\label{fig:life}
\end{figure}

\begin{itemize}
  \item \textbf{Alive.} The rover infiltrates (Section~\ref{sec:rover}); heat footprints are laid under its wheels and the parked hull. Damage slows it and, below 60 HP, makes it smoke.
  \item \textbf{Wreck.} \code{destroyIntruder()} counts a kill, stops the rover, tells the squad to ignore the wreck, saves every part's material and replaces it by a charred one
  (albedo $\times0.1+(0.015,0.013,0.012)$, roughness 0.95, metalness 0, heat 0.95, no emission), turns all parts except the hull core
  (\emph{Chassis, Upper hull, Motor + battery box}) into rigid-body debris and adds a scorch mark.
  \item \textbf{Burning.} For 15\,s the wreck emits fire, smoke and embers in proportion to the remaining burn fraction;
  the hull cools as $h=0.3+0.65\,e^{-t/40\,\text{s}}$ (visible in the thermal view) and deposits heat on the
  ground (0.25 per second). After 25\,s the wreck is removed and 3\,s later the slot is free; T107 checks wreck $\to$ gone $\to$ alive.
\end{itemize}

\subsection{Engagement rules}\label{sec:engage}
The UAVs' target selection, lock times, friend-or-foe and keep-out rules, wreck discrimination and the combined strike
belong to the squad (Section~\ref{sec:squad}); the thermal detection behind them is Chapter~\ref{sec:vision}. Manual flight
disables the autonomous launch, and \texttt{Enter} releases a rail round at the camera's aim point.

\subsection{Projectiles}
\sheetfig{s16_weapons}
\paragraph{Proportional-navigation homing.} A UAV missile (speed profile in Section~\ref{sec:squad}) homes on its laser spot, the ground point at the centre of the gimbal image. With the spot at
$\vect r$ (vector from the missile), $\hat{\vect v}$ its heading, $V$ its speed and $\hat{\vect l}=\vect r/|\vect r|$ the line of sight,
the commanded turn rate is
\begin{equation}
  \boldsymbol\omega=\frac{4V}{|\vect r|^2}\,(\hat{\vect v}\times\vect r)+2\,(\hat{\vect v}\times\hat{\vect l}),\qquad |\boldsymbol\omega|\le6\ \text{rad/s},
\end{equation}
and the heading is rotated about $\hat{\boldsymbol\omega}$ by $|\boldsymbol\omega|\Delta t$ (Rodrigues' formula). The first term is proportional navigation with $N=4$ for a
stationary spot, the second a heading-error term. The turn limit models limited manoeuvrability. The proximity fuse radius is 0.5\,m, a closest-approach
fuse fires within 3\,m of the spot, and a round self-destructs after 6\,s (T100: a round launched
$30^\circ$ off the line of sight detonates 0.50\,m from the spot after 0.32\,s). Substeps are at most $1/120$\,s.

\paragraph{Ballistic shell.} The tank shell flies at 120\,m/s under gravity without drag; T101 finds a range of 203.8\,m against
$v^2\sin2\alpha/g$ (within $-0.24$\,\%).

\paragraph{Tracer and segment tests.} The AA tracer flies at 1000\,m/s for at most 2.5\,s with no blast at burn-out. At 1000\,m/s a round moves 16.7\,m per
$1/60$\,s step, so impacts are tested as \emph{segments} from the previous to the new position against every target sphere (radius plus fuse), every static solid box
(slab test) and the ground, taking the earliest entry $t\in[0,1]$; nothing tunnels (T102 finds the small target). Blast power: missile 1.0, shell 1.6, tracer 0.15.

\subsection{Blast damage}
\sheetfig{s16_effects}
A blast of power $p$ damages every living intruder within $R=1.6\,p$ metres (\code{kDamageRadiusPerPower}) and the one hit directly. An armoured rover has
\code{kRoverHealth}\,$=100$ HP and loses $30\,p\,(1-\tfrac12\min(d/R,1))$ per blast ($d=0$ for a direct hit), so four direct power-1 hits destroy it (T105: five hits to kill, nothing at $3\,R$);
earlier hits only raise an \code{EnemyHit} event, sparks and dust. Controllable objects within $6\,p$ metres receive \code{onBlast()}; the tank and the AA gun scorch (darken their camouflage). A blast with $p\ge0.3$
spawns a flash (0.12\,s), a fireball of 10--25 fire puffs, a smoke column of 15--40 puffs with a 6--12\,s life, 30--60 sparks, 8--15 embers, a ground shock-wave ring, 10--20 dust
puffs, a decaying point light, camera shake within 25\,m (decaying over 0.8\,s) and a scorch mark of radius $0.9\,p$ (cap 96 marks, 120\,s life). T103 checks that all seven main
particle types appear, the light decays 1 $\to$ 0 and one scorch is left.

\subsection{Rigid-body debris}
Each non-core part becomes an independent rigid body. At the moment of the blast the part's world matrix $W_0$ and origin $\vect c_0$ are stored. The initial velocity points
away from the blast with a lift,
$\vect v_0=\mathrm{normalize}(\vect a+\text{jitter})\,(2.5+3.5\,r)\sqrt p+(0,1.5+3r',0)$ with $\vect a=\vect c_0-\vect b$, its $y$ component clamped to $\ge0$ and raised by 0.6\,m, and the spin has a pseudo-random axis and a rate of
4--14\,rad/s. Each substep ($\le1/120$\,s) applies gravity, moves the body, and composes the rotation with a Rodrigues rotation about the spin axis. The world matrix is
\begin{equation}
  W(t)=T(\vect c(t))\,R(t)\,T(-\vect c_0)\,W_0 ,
\end{equation}
i.e.\ the part is rotated about its own origin and moved to the integrated position. On ground contact (terrain $+0.03$\,m): restitution 0.3 in the normal direction,
horizontal speed $\times0.7$ and spin $\times0.6$ per contact; a piece comes to rest below 0.25\,m/s. T106: 23 pieces, peak height 3.25\,m, all at rest on top of the terrain within
8\,s. Because the hull core stays and the other 23 parts fly, a wreck keeps a recognisable chassis.

\begin{figure}[h]
  \centering
  \begin{subfigure}{0.325\linewidth}\picfour{thermal_fpv_strike}{\linewidth}\caption{Thermal FPV at the strike}\end{subfigure}
  \begin{subfigure}{0.325\linewidth}\picfour{uav_missile_launch}{\linewidth}\caption{Missile launch}\end{subfigure}
  \begin{subfigure}{0.325\linewidth}\picfour{explosion_early}{\linewidth}\caption{Explosion, +0.15\,s}\end{subfigure}\\[1mm]
  \begin{subfigure}{0.325\linewidth}\picfour{explosion_mid}{\linewidth}\caption{Explosion, +1\,s}\end{subfigure}
  \begin{subfigure}{0.325\linewidth}\picfour{explosion_smoke}{\linewidth}\caption{Smoke column, +5\,s}\end{subfigure}
  \begin{subfigure}{0.325\linewidth}\picfour{wreck_burning}{\linewidth}\caption{Burning wreck, +9\,s}\end{subfigure}
  \caption{The strike sequence (captures at the delays set in \code{Capture.cpp}).}
  \label{fig:strike}
\end{figure}

